% Part: history
% Chapter: biographies
% Section: gerhard-gentzen

\documentclass[../../../include/open-logic-section]{subfiles}

\begin{document}

\olfileid{his}{bio}{gen}

\olsection{گیرہارٹ گینٹسن}

\olphoto{gentzen-gerhard}{گیرہارٹ گینٹسن}

گیرہارٹ گینٹسن دی پہچان خاص کر کے ساختی نظریۂ ثبوت دے بانی وجوں اے۔ خاص طور اُتے اوہنے قدرتی اخذ تے سیکوئنٹ کیلکولس دے !!{derivation} نظام بنائے۔ اوہ 24 نومبر 1909 نوں جرمنی دے گرائفسوالڈ وچ جمیا۔ گیرہارٹ نے تیاری والے سکول جان توں پہلاں تِن سال گھر وچ پڑھائی کیتی۔ سکول وچ اوہ پڑھائی پکھوں اپنے بہتے جماعتیاں توں پچھے سی۔ ایس دے باوجود اوہ بڑا ہونہار طالب علم سی تے ریاضی وچ اوہدی خاص قابلیت دِسدی سی۔ اوہ کئی طرح دے کم پسند کردا سی، مثلاً اپنی ماں لئی نظماں تے سکول دے تھیٹر لئی ڈرامے وی لکھدا سی۔

گینٹسن نے یونیورسٹی دی پڑھائی گرائفسوالڈ یونیورسٹی توں شروع کیتی، پر فیر گوٹنگن، میونخ تے برلن وی گیا۔ اوہنے 1933 وچ گوٹنگن یونیورسٹی توں ہرمان وائل دی نگرانی تھلّے ڈاکٹریٹ لئی۔ (اوہدے کم دے بہتے حصے دی نگرانی پاؤل بیرنائز نے کیتی، پر نازیاں نے بیرنائز نوں یونیورسٹی توں کڈھ دتا سی۔) 1934 وچ گینٹسن نے ڈیوڈ ہلبرٹ دے معاون وجوں کم شروع کیتا۔ اوسے سال اوہنے اپنے مقالیاں \emph{Untersuchungen \"{u}ber das logische Schlie\ss en I--II
  [منطقی اخذ دی پڑتال، حصے ۱ تے ۲]} وچ سیکوئنٹ کیلکولس تے قدرتی اخذ دے !!{derivation} نظام بنائے۔ اوہنے 1936 وچ پیانو دے مسلّمیاں دا بے ٹکراؤ ہونا ثابت کیتا۔

گینٹسن دا نازیاں نال تعلق گنجل دار اے۔ جیہڑے ویلے اوہدے استاد بیرنائز نوں جرمنی چھڈن لئی مجبور کیتا گیا، اوسے ویلے گینٹسن نازی نیم فوجی تنظیم ایس اے دی یونیورسٹی شاخ وچ شامل ہویا۔ کئی ہور جرمن لوکاں وانگ اوہ نازی پارٹی دا رکن سی۔ جنگ دوران اوہ فضائی انٹیلیجنس یونٹ وچ مواصلات دا افسر سی۔ پر 1942 وچ شدید ذہنی پریشانی کارن اوہنوں ڈیوٹی توں فارغ کر دتا گیا۔ ایہ صاف نہیں کہ گینٹسن نازی پارٹی دا وفادار سی یا اپنی علمی ترقی پکی کرن لئی پارٹی وچ شامل ہویا۔

1943 وچ پراگ دی جرمن یونیورسٹی دے ریاضی ادارے وچ گینٹسن نوں اک علمی عہدے دی پیشکش ہوئی، جیہڑی اوہنے منّ لئی۔ پر 1945 وچ پراگ دے شہری جرمن قبضے دے خلاف اُٹھ کھلوتے۔ سوویت فوج شہر وچ آئی تے یونیورسٹی دے سارے پروفیسراں نوں گرفتار کر لیا۔ نازی تنظیماں دی رکنیت کارن گینٹسن نوں جبری مزدوری دے کیمپ لے گئے۔ اوہ 4 اگست 1945 نوں 35 سال دی عمر وچ اپنی کوٹھڑی وچ خوراک دی کمی نال مر گیا۔

\begin{reading}
گینٹسن دی پوری سوانح عمری لئی \citet{Menzler-Trott2007} ویکھو۔ نازی راج تھلّے ریاضی داناں بارے اک دلچسپ لکھت \citet{Segal2014} دی اے، جیہدے وچ گینٹسن دی زندگی دا مختصر ذکر وی اے۔ منطقی اخذ اُتے گینٹسن دے مقالے اصل جرمن وچ ملدے نیں \citep{Gentzen1935a,Gentzen1935b}۔ \citet{Gentzen1969} وچ گینٹسن دے مقالیاں دے انگریزی ترجمے اکّو جلد وچ اکٹھے کیتے گئے نیں تے اوہدے وچ اوہدی زندگی دا مختصر خاکہ وی اے۔
\end{reading}

\end{document}
